% ============================================================================
% Methodology
%
% Purpose
% -------
% Explain HOW the research was conducted, not WHAT was discovered.
%
% Reader should understand:
%   • the overall research workflow
%   • why the analysis is reproducible
%   • the research dataset
%   • variables analysed
%   • statistical methods
%   • machine learning models
%   • independent R verification
%   • automated reporting pipeline
%
% Do NOT discuss:
%   • model performance
%   • feature importance
%   • educational interpretation
%   • conclusions
%
% Transition:
% This section prepares the reader to understand and trust the Results.
% ============================================================================


\section{Methodology} 

This study adopted a reproducible analytical workflow to investigate factors associated with successful transition into first-year university mathematics. 
The workflow was designed to ensure that each stage of the analysis, from data preparation through statistical modelling, independent verification, and report 
generation, could be reproduced from the original research dataset using an automated build process. This approach emphasises transparency, repeatability, and 
traceability, allowing both analytical results and publication outputs to be generated consistently from the same underlying evidence.

% --------------------------------------------------------------------------
% Research workflow
%
% Explain the overall analytical process before introducing individual
% components such as the dataset and modelling methods.
%
% This subsection establishes reproducibility as the guiding principle.
% --------------------------------------------------------------------------
\subsection{Research Workflow}

This study followed a reproducible analytical workflow consisting of data preparation, statistical analysis, machine learning model development, independent 
verification using R, and automated generation of publication-ready tables, figures, and manuscript components. Each stage formed part of an integrated build 
process, ensuring that analytical outputs could be reproduced directly from the underlying research dataset.

Figure~\ref{fig:workflow} summarises the reproducible analytical workflow used in this study, from preparation of the approved research dataset through 
statistical analysis, machine learning, independent verification, and automated publication outputs.

\begin{figure}[htbp]
\centering

\fbox{
\begin{minipage}{0.92\textwidth}
\centering

\textbf{Institutional Source Data}

\vspace{0.15cm}
$\downarrow$

\vspace{0.15cm}
\textbf{Data Preparation}

\vspace{0.15cm}
$\downarrow$

\vspace{0.15cm}
\textbf{Research Dataset}

\vspace{0.15cm}
$\downarrow$

\vspace{0.15cm}
\begin{tabular}{c@{\hspace{1.5cm}}c}
\fbox{\parbox{0.30\textwidth}{\centering
\textbf{Statistical Analysis}
}}
&
\fbox{\parbox{0.30\textwidth}{\centering
\textbf{Machine Learning Analysis}
}}
\end{tabular}

\vspace{0.15cm}

\[
\searrow \hspace{2.8cm} \swarrow
\]

\vspace{0.15cm}
$\downarrow$

\vspace{0.15cm}
\textbf{Independent Verification (R)}

\vspace{0.15cm}
$\downarrow$

\vspace{0.15cm}
\textbf{Tables • Figures • Generated LaTeX Values}

\vspace{0.15cm}
$\downarrow$

\vspace{0.15cm}
\textbf{Publication-ready Paper}

\end{minipage}
}

\caption{Reproducible analytical workflow implemented by the Transition Analysis Toolkit. Institutional data are transformed into an approved 
de-identified research dataset before parallel statistical and machine learning analyses are independently verified in R. Automated generation 
of tables, figures, LaTeX variables, and manuscript content ensures consistency between the analytical outputs and the final publication.}

\label{fig:workflow}

\end{figure}



% --------------------------------------------------------------------------
% Research Dataset
%
% Purpose
% -------
% Describe the research data used in this study.
%
% Reader should understand:
%   • where the data came from
%   • who is represented
%   • ethical considerations
%   • why the dataset is appropriate
%
% Do NOT discuss:
%   • model results
%   • predictor importance
%   • interpretation
%
% Transition:
% The next subsection describes the variables derived from this dataset.
% --------------------------------------------------------------------------
\subsection{Research Dataset} 

The study used de-identified institutional data from students enrolled in a first-stage undergraduate mathematics subject offered during the Autumn 2026 teaching session. The dataset was prepared specifically for research purposes using an approved de-identification process that removed all personally identifying information while preserving the academic and demographic characteristics required for analysis. The resulting dataset provided the foundation for all subsequent statistical analyses, machine learning models, independent verification procedures, and automated reporting undertaken in this study.




\subsection{Variables} 

The analytical dataset comprised variables describing students' prior academic preparation, previous university mathematics experience, demographic characteristics, and current academic outcomes. Variables were selected to reflect factors commonly associated with transition into first-year university mathematics and were derived from institutional records following the approved research data preparation process. A complete description of each variable is provided in Table~\ref{tab:variables}.

\begin{table}[ht]
\centering
\caption{Summary of variables included in the analytical dataset.}
\label{tab:variables}

\begin{tabular}{ll}
\hline
Variable & Description \\
\hline
\multicolumn{2}{c}{\emph{Placeholder for variable summary table}} \\
\hline
\end{tabular}
\end{table}



\subsection{Statistical Analysis} 

Statistical analyses were undertaken to summarise the characteristics of the research dataset and to examine relationships between student background variables and academic outcomes. Descriptive statistics were generated to provide an overview of the cohort, while inferential analyses supported comparison of predictor variables prior to the development of machine learning models. These analyses established a statistical baseline against which the predictive models could subsequently be interpreted.




\subsection{Machine Learning Models} 

Machine learning models were developed to investigate the relative importance of factors associated with successful transition into first-year university mathematics. Decision tree and random forest classifiers were selected because they provide both competitive predictive performance and interpretable measures of feature importance. Hyperparameter tuning and model evaluation were undertaken using reproducible workflows to ensure consistent comparison across modelling approaches.



\subsection{Independent Verification} 

To strengthen confidence in the analytical results, an independent verification pipeline was implemented using the R statistical computing environment. Key statistical summaries and model outputs were reproduced independently of the primary Python implementation to confirm the consistency and reliability of the analytical workflow. This independent verification formed an additional quality assurance stage prior to the generation of publication outputs.



\subsection{Automated Reporting} 

The final stage of the analytical workflow automatically generated publication-ready tables, figures, model summaries, and manuscript components directly from the verified analytical outputs. By integrating analysis and reporting within a single reproducible workflow, the potential for manual transcription errors was reduced while ensuring that all reported results remained synchronised with the underlying research data and analyses.



The methodology described above established a transparent and reproducible analytical framework for investigating factors associated with successful transition into first-year university mathematics. The following section presents the outcomes of this workflow, beginning with an evaluation of predictive model performance before examining the relative importance and educational interpretation of the identified predictors.


